% ---------------------------------------------------------------------
%  Chapter 20 : Validation --- Independent Simulators & Cumulant Estimators
% ---------------------------------------------------------------------
\chapter{Validation: Independent Simulators and Cumulant Estimators}\label{ch:simulators}

Everything in the preceding chapters builds \emph{one} half of the repository: the
analytic pipeline that takes an \msrjd{} action, enumerates Feynman diagrams,
assigns propagators and vertices, and computes connected correlation functions
(cumulants) order by order in the loop expansion. This chapter documents the
\emph{other} half --- the half whose entire job is to not believe the first one.

The files described here are a set of completely independent numerical
simulators that solve the \emph{same} stochastic equations of motion directly,
by brute-force time-stepping, plus a set of estimators that turn the raw
simulated time series into empirical estimates of the very same cumulant slices
the pipeline predicts. The two are then plotted on top of each other. If the
diagrammatic machinery is correct, the curves overlap, within Monte-Carlo error
bars. This is the \term{trust chain}: a closed loop in which an analytic
prediction is checked against a numerical experiment that shares \emph{no code}
with the prediction. You met one link of that chain already, at the end of
Chapter~\ref{ch:quickstart}, where the two-loop value $C_{xx}(0) = 0.9496$ was
confirmed by a simulation that measured $0.9464 \pm 0.0009$. This chapter is the
full story behind that one line of output.

\begin{note}
\textbf{These files are not part of the pipeline.} They never run during a normal
cumulant computation. They exist solely so that a skeptical reader can confirm
the diagrams give the right answer. The repository's confidence in every analytic
result rests on these brute-force checks agreeing. When you read ``the pipeline''
elsewhere in this manual, the simulators are \emph{not} included; they are the
external referee.
\end{note}

The two primary files are
\file{models/ou\_langevin\_sim\_numba.py} (the Euler--Maruyama SDE integrators)
and \file{models/cumulant\_estimator.py} (the $k$-point connected-cumulant
estimator). A wider supporting cast --- the Hawkes point-process simulators, the
spatial-field simulators, and the coupled reaction--diffusion oracle --- is
characterised at the end of the chapter. We open every tool from the ground up:
unlike the pipeline, this subsystem touches no SageMath, no \code{nauty}, no
\code{sympy}. It is deliberately lean, and that cleanness is the whole point ---
the validators share no symbolic machinery with the code they validate.

\section{Why a separate validation layer?}

A number that comes out of a diagrammatic calculation is the end of a long chain
of symbolic manipulations: a Taylor expansion of the action, a graph enumeration,
an automorphism-counting symmetry factor, a multidimensional loop integral. A bug
anywhere in that chain produces a wrong number that still \emph{looks} like a
correlation function --- it is smooth, it decays, it is symmetric in $\tau$. There
is nothing in the output itself that flags the error. The only way to catch such
a bug is to compute the same quantity a \emph{second}, completely different way,
and compare.

The cheapest fully independent second computation for a stochastic differential
equation is to simply \emph{simulate} it: discretise time, take small steps,
add the right amount of random noise each step, and record the trajectory. Run it
long enough and the time-averaged statistics of that trajectory converge (by
ergodicity) to the stationary correlation functions the theory predicts. This
shares nothing with the diagrammatics --- there is no action, no graph, no
integral. A human has hand-coded the equation of motion as ordinary numeric
arithmetic. If both routes land on the same number, the probability that they
both contain a compensating bug is negligible, and the analytic result is
licensed.

The phrase ``within error bars'' is load-bearing. A simulation is a Monte-Carlo
estimate: it has statistical noise that shrinks like $1/\sqrt{N}$ in the number of
independent runs (or the length of one run), plus a few systematic biases from the
finite step size and the finite bin width that we will quantify below. The
comparison is therefore never an exact equality; it is ``the theory value sits
inside the simulation's confidence interval.'' When the example notebook reports
$0.9464 \pm 0.0009$ against a theory value of $0.9496$, the residual $0.003$ gap
is about three standard errors --- accounted for by the two-loop truncation plus
the estimator's finite-bin and finite-$T$ biases.

\section{The two simulator families}

There are two physically distinct kinds of dynamics the repository validates,
and correspondingly two families of simulator.

\begin{description}
  \item[Langevin / SDE simulators] integrate a stochastic \emph{differential}
        equation, $\dot x = f(x) + \text{noise}$, where the state $x$ is a
        continuous real variable (or field). These cover the OU-type theories
        (\file{ou\_langevin\_sim\_numba.py}), the spatial scalar fields
        (\file{spatial\_field\_*\_sim.py}), and the coupled reaction--diffusion
        systems (\file{coupled\_rd\_1d\_sim.py}). The workhorse is the
        Euler--Maruyama scheme; stiffer linear parts use an exact
        Ornstein--Uhlenbeck step, a spectral exponential-Euler step (ETD1), or a
        semi-implicit Crank--Nicolson step.
  \item[Point-process / Hawkes simulators] integrate a self-exciting Poisson
        point process (\file{hawkes\_sim\_*numba.py}). The state is a discrete
        spike count: each timestep draws $\mathrm{Poisson}(\lambda_i\,\mathrm{d}t)$
        spikes and feeds them back into the rate. This is a discrete-event
        process, not an SDE, and it is what makes the \emph{factorial-cumulant}
        machinery of \S\ref{sec:factorial} necessary.
\end{description}

The same estimator file consumes both: a simulator returns a binned trajectory
of shape \code{(npop, n\_bins)}, and the estimator cross-correlates that to
produce a cumulant slice. Which leg is treated as a smooth continuous field and
which as a discrete spike train is selected by a per-leg \emph{field-type} label
(\S\ref{sec:fieldtypes}).

\subsection{Where it sits in the end-to-end pipeline}

The wiring is a diamond: one theory file feeds \emph{both} the pipeline and the
simulator, and the two output streams meet in a comparison notebook.

\begin{center}
\begin{tikzpicture}[
    box/.style={draw,rounded corners,align=center,font=\small,inner sep=5pt},
    >=Stealth, node distance=7mm]
  \node[box] (theory) {theory file\\ (action $S$, params)};
  \node[box, right=28mm of theory] (pipe) {\textsc{pipeline}\\ enumerate\\$\to$ assign $\to$ integrate};
  \node[box, below=16mm of pipe] (sim) {\textsc{simulator}\\ Euler--Maruyama\\ / Poisson};
  \node[box, right=24mm of pipe] (res) {\code{res['C\_tau']}\\ analytic cumulants};
  \node[box, right=24mm of sim] (est) {\textsc{estimator}\\ sim cumulant\\ slices};
  \node[box, right=22mm of res, yshift=-8mm] (cmp) {sim-compare\\ notebook\\ (overlay plot)};
  \draw[->] (theory) -- (pipe);
  \draw[->] (theory.south) |- (sim.west);
  \draw[->] (pipe) -- (res);
  \draw[->] (sim) -- (est);
  \draw[->] (res.east) -- (cmp.north west);
  \draw[->] (est.east) -- (cmp.south west);
\end{tikzpicture}
\end{center}

Three facts about that diagram matter for the rest of the chapter:

\begin{itemize}
  \item \textbf{The simulator is handed the pipeline's resolved parameters.} In
        the sim-compare notebooks the call is literally
        \code{fp = res['\_resolved']['parameters']} followed by
        \code{mu = float(fp['mu'])}, and so on. The simulator gets the
        \emph{same} numbers the theory used, cast to plain Python floats. It does
        \emph{not} read the action --- a human has hand-coded each equation of
        motion to match a specific theory file, and the docstring of each
        \code{sim\_*} function names the matching \file{theories/*.theory.py}.
  \item \textbf{The simulator output feeds the estimator.} A simulator returns a
        binned trajectory array of shape \code{(npop, n\_bins)}; the estimator
        cross-correlates that to produce a cumulant slice \code{(n\_tau,)} or a
        stack \code{(k-1, n\_tau)}.
  \item \textbf{The estimator output feeds the comparison plot.} The notebooks
        overlay the simulation's \code{sim['C']} against the theory's
        \code{res['C\_tau\_by\_ell']} (per loop order) or \code{res['C\_tau']}
        (loop-summed). \emph{That overlap is the trust-chain verdict.}
\end{itemize}

\section{Euler--Maruyama from scratch}

The scalar simulators integrate an It\^o stochastic differential equation
\begin{equation}
  \dot x \;=\; f(x) \;+\; \sqrt{2D}\,\eta(t),
  \qquad
  \avg{\eta(t)\,\eta(t')} \;=\; \delta(t-t'),
  \label{eq:sde}
\end{equation}
where $f(x)$ is the deterministic \term{drift} and $\eta$ is unit-strength
Gaussian white noise. For the flagship example --- the quartic double well in
\code{sim\_ou\_quartic\_numba} --- the drift is
\[
  f(x) \;=\; -\mu x \;-\; \varepsilon x^{3},
\]
so the equation of motion is $\dot x = -\mu x - \varepsilon x^{3} + \sqrt{2D}\,
\eta$. This is exactly the \msrjd{} action quoted in the docstring
(\file{ou\_langevin\_sim\_numba.py:29}),
\[
  S \;=\; \int \dd t\; \tilde x\,\bigl((\Dt + \mu)\,x + \varepsilon x^{3}\bigr)
        \;-\; D\,\tilde x^{2},
\]
where $\tilde x$ (written \code{xt} in the theory text) is the \msrjd{} response
field. The $-D\,\tilde x^{2}$ term is the noise vertex; its coefficient $D$ fixes
the noise strength $\avg{\xi\,\xi} = 2D\,\delta(\tau)$ with $\xi = \sqrt{2D}\,\eta$.

\subsection{The defining feature: $\sqrt{\dd t}$ noise scaling}

\term{Euler--Maruyama} is the stochastic analogue of the forward-Euler scheme for
ordinary differential equations. Over a step $\dd t$ the deterministic part
advances by $\dd t\,f(x)$, exactly as in a deterministic Euler step. The new
ingredient is the integrated noise increment $\int_t^{t+\dd t}\sqrt{2D}\,\eta\,
\dd t'$, which is a Gaussian with mean $0$ and variance $2D\,\dd t$. (The variance
follows from $\mathrm{Var}\!\left[\int \eta\right] = \int\!\!\int \delta = \dd t$:
the double integral of the delta-correlation over the step collapses to the step
length.) Therefore one Euler--Maruyama step is
\begin{equation}
  x(t+\dd t) \;=\; x(t) \;+\; \dd t\, f(x) \;+\; \sqrt{2D\,\dd t}\;N(0,1).
  \label{eq:em}
\end{equation}
The crucial detail is the $\sqrt{\dd t}$ scaling of the noise term --- not $\dd t$.
A deterministic Euler step would scale \emph{everything} by $\dd t$; the noise
must scale by $\sqrt{\dd t}$ because variance is what adds linearly in time, so
the standard deviation grows like $\sqrt{\dd t}$. This single fact is the whole
difference between Euler and Euler--Maruyama.

In code, the noise prefactor is precomputed once outside the loop
(\file{ou\_langevin\_sim\_numba.py:74}) and the step itself is three lines
(\file{ou\_langevin\_sim\_numba.py:80--82}):

\begin{lstlisting}
sqrt_2D_dt = np.sqrt(2.0 * D * dt_sim)
...
x = (x
     + dt_sim * (-mu * x - eps * x * x * x)
     + sqrt_2D_dt * np.random.randn())
\end{lstlisting}

\code{np.random.randn()} draws one $N(0,1)$ sample; multiplying by
\code{sqrt\_2D\_dt} turns it into the correctly-scaled noise increment. Note the
cubic is written \code{eps * x * x * x} rather than \code{eps * x**3}: repeated
multiplication is cheaper than a power call inside a hot loop, and (as we will see)
this code runs millions of times per simulation.

\subsection{Weak order one and the $\dd t \ll 1/\mu$ rule}

Euler--Maruyama is a \term{weak order-1.0} scheme. ``Weak'' means we care about the
accuracy of \emph{averages} and \emph{moments} of $x$ (which is what a correlation
function is), not about the pathwise accuracy of any single trajectory. ``Order
$1.0$'' means the bias in those stationary statistics scales as $O(\dd t)$ ---
halve the step, halve the bias. The practical consequence, stated in the docstring
(\file{ou\_langevin\_sim\_numba.py:36}), is that the step must be small relative to
the relaxation time $1/\mu$: $\dd t_{\mathrm{sim}} \ll 1/\mu$. If you step too
coarsely, the simulation systematically misrepresents the fast relaxation of $x$
and biases every measured statistic.

\begin{defn}
\textbf{Stationary variance sanity check.} For the \emph{pure} OU process
($\varepsilon = 0$, $f = -\mu x$) the stationary distribution is Gaussian with
variance $D/\mu$, so the equal-time correlation is $C_{xx}(0) = \avg{x^{2}} =
D/\mu$. This is the analytic target at tree level; a finite $\varepsilon$ shifts it
by $O(\varepsilon D^{2}/\mu^{3})$ (\file{ou\_langevin\_sim\_numba.py:64--66}). It is
exactly the number the pipeline's loop expansion should reproduce, and a quick way
to confirm a simulation is behaving: set $\varepsilon = 0$ and check the measured
variance lands on $D/\mu$.
\end{defn}

The sign conventions follow the theory file
(\file{ou\_langevin\_sim\_numba.py:60--63}): positive $\mu$ gives a stable origin
(sub-critical regime), $\mu = 0$ is a pitchfork bifurcation, and $\mu < 0$ is a
genuine double well with minima at $x = \pm\sqrt{|\mu|/\varepsilon}$. The
equilibrium Boltzmann potential is $U(x) = \tfrac{\mu}{2}x^{2} +
\tfrac{\varepsilon}{4}x^{4}$ (and $+\tfrac{\gamma}{6}x^{6}$ for the sextic
variant), bounded below --- hence stable and normalisable --- when all coefficients
are positive.

\subsection{The bin accumulator}

All five functions in \file{ou\_langevin\_sim\_numba.py} share one skeleton: seed
the random generator, initialise the state, loop \code{n\_steps} doing an Euler
step plus a bin accumulation, and emit a bin-averaged trajectory. The simulation
runs at a fine step \code{dt\_sim} but records output only every
\code{bin\_size\_steps} steps, averaging $x$ over that window into one output bin
(\file{ou\_langevin\_sim\_numba.py:84--89}):

\begin{lstlisting}
steps_in_bin += 1
if steps_in_bin >= bin_size_steps:
    if cur_bin < n_bins:
        x_bins[0, cur_bin] = accum / bin_size_steps
        accum = 0.0
    cur_bin += 1
    steps_in_bin = 0
\end{lstlisting}

The guard \code{if cur\_bin < n\_bins} stops writing once the output buffer is
full; any extra Euler steps beyond \code{bin\_size\_steps}\,$\times$\,\code{n\_bins}
are simply ignored. The output shape is always \code{(npop, n\_bins)}, matching
the Hawkes \code{voltage\_bins} convention, so the estimator can consume an OU
trajectory with \code{field\_types=['dv', ...]} exactly as it consumes a Hawkes
voltage trace.

\begin{gotcha}
\textbf{Bin finely, or bias $C_{xx}(0)$ low.} A coarse bin (one whose width
\code{dt\_bin} is not $\ll 1/\mu$) averages the field over the bin and smooths away
the fastest fluctuations, so the bin-averaged variance \emph{underestimates} the
instantaneous $C_{xx}(0)$ the theory reports. The example notebook sets
\code{dt\_bin}\,$= 0.02$ for $\mu \sim 1$ precisely to keep this bias below the
statistical error. This is a separate requirement from the $\dd t_{\mathrm{sim}}
\ll 1/\mu$ rule: one governs the integration step, the other the recording window.
\end{gotcha}

\section{Beyond plain Euler: removing bias from the linear part}

Plain Euler--Maruyama carries an $O(\dd t)$ bias in \emph{everything}, including
the linear relaxation. Three of the simulator families improve on this by treating
the \emph{linear} part of the dynamics exactly, leaving only the genuinely
nonlinear pieces (and, for colored noise, the $x$ step) at $O(\dd t)$. Each trick
removes a specific bias.

\subsection{Exact Ornstein--Uhlenbeck discretisation (colored noise)}
\label{sec:colored}

Some theories use temporally correlated (\term{colored}, Lorentzian) noise rather
than white noise. \code{sim\_ou\_quartic\_colored\_numba} realises this by running
a \emph{second} SDE --- an auxiliary Ornstein--Uhlenbeck process $\xi$ --- and using
its output as the forcing of $x$:
\begin{align}
  \dot\xi &= -\xi/\tau_c + (\sqrt{2D}/\tau_c)\,\eta(t)
            &&\Longrightarrow&
            \avg{\xi(t)\,\xi(t')} &= (2D/\tau_c)\,\ee^{-|t-t'|/\tau_c}, \\
  \dot x  &= -\mu x - \varepsilon x^{3} + \xi(t). &&&&
\end{align}
The key trick is that a \emph{linear} OU SDE has a closed-form one-step update ---
it does not need Euler--Maruyama at all. The exact update is
\begin{equation}
  \xi(t+\dd t) \;=\; \text{decay}\cdot\xi(t) \;+\; \sigma\,N(0,1),
  \qquad
  \text{decay} = \ee^{-\dd t/\tau_c},
  \qquad
  \sigma^{2} = (2D/\tau_c)(1 - \text{decay}^{2}).
  \label{eq:exactou}
\end{equation}
This preserves the stationary variance for \emph{any} ratio $\dd t/\tau_c$: there
is no $O(\dd t)$ bias in the noise statistics at all, only in the $x$ Euler step.
In code (\file{ou\_langevin\_sim\_numba.py:146--148}):

\begin{lstlisting}
decay = np.exp(-dt_sim / tauc)
var_factor = 1.0 - decay * decay
sigma_xi = np.sqrt((2.0 * D / tauc) * var_factor)
\end{lstlisting}

A subtle implementation detail: at each step the code draws the new noise
\code{xi\_new} but steps $x$ with the \emph{previous}-step \code{xi}, then commits
both (\file{ou\_langevin\_sim\_numba.py:160--166}):

\begin{lstlisting}
xi_new = decay * xi + sigma_xi * eta
drift = -mu * x - eps * x * x * x
x_new = x + dt_sim * (drift + xi)   # NB: previous-step xi, not xi_new
x = x_new
xi = xi_new
\end{lstlisting}

\begin{gotcha}
\textbf{Colored-noise factor of 2.} In the white limit $\tau_c \to 0$ this colored
driver tends to $\avg{\xi\xi} \to 4D\,\delta(\tau)$ --- note \textbf{4D}, not 2D ---
because of the $2D/\tau_c$ coefficient convention. To compare the colored
simulator against the white simulator you must use $D_{\mathrm{white}} = 2\,
D_{\mathrm{colored}}$ (equivalently, halve $D$ when crossing over). The docstring
spells this out (\file{ou\_langevin\_sim\_numba.py:127--132}); forgetting it
produces a clean-looking curve that is wrong by a constant factor.
\end{gotcha}

\subsection{Spectral exponential-Euler (ETD1) for stiff spatial fields}

The spatial simulators (\file{spatial\_field\_1d\_sim.py} and its 2D/3D siblings)
integrate a \emph{field} SDE on a periodic ring or torus, e.g.\ the
non-conserved $\phi^{4}$ dynamics
\[
  \Dt \phi(x,t) \;=\; -\mu\phi + D\,\partial_x^{2}\phi - \lambda\phi^{3} + \eta(x,t),
  \qquad
  \avg{\eta(x,t)\,\eta(x',t')} = 2T\,\delta(x-x')\,\delta(t-t').
\]
A naive Euler--Maruyama in real space is unstable here because the diffusion
operator $D\,\partial_x^{2}$ makes the system \term{stiff}: its largest eigenvalue
grows like $D\,k_{\max}^{2}$, and an explicit step must resolve that fastest mode
or blow up. The fix is \term{ETD1} (exponential time-differencing, order 1),
working in Fourier space:

\begin{itemize}
  \item Each Fourier mode $\hat\phi_k$ obeys a \emph{linear} OU SDE with rate
        $\omega_k = \mu + D k^{2}$ (the lattice version is $\omega_k = \mu +
        (2D/\dd x^{2})(1 - \cos(k\,\dd x))$), which is integrated \emph{exactly}
        with the same OU closed form as \eqref{eq:exactou}. Consequently the
        $\lambda = 0$ stationary spectrum is unbiased in $\dd t$.
  \item The nonlinear forcing $-\lambda\phi^{3}$ is added through the ETD1
        integrating factor $(1 - \ee^{-\omega\,\dd t})/\omega$
        (\file{spatial\_field\_1d\_sim.py:77}).
\end{itemize}

The per-mode stationary variance is $\avg{|\hat\phi_k|^{2}} = T N^{2}/(L\,\omega_k)$
(\file{spatial\_field\_1d\_sim.py:78}), and the per-step OU increment carries
variance $(1 - \ee^{-2\omega\,\dd t})$ times that
(\file{spatial\_field\_1d\_sim.py:79}). Real-mode versus interior-complex-mode
bookkeeping (\file{spatial\_field\_1d\_sim.py:83--87} and~\code{:114--118})
handles the Hermitian structure of the real FFT. Derivative vertices --- Model~B's
conserved $g_{\mathrm{lap}}\,\partial_x^{2}(\phi^{2})$, the Burgers
$-(\lambda/2)\partial_x(\phi^{2})$, the KPZ $+(\lambda/2)(\partial_x\phi)^{2}$ ---
enter through spectral multipliers \code{lap\_eig} and \code{ik\_eig}
(\file{spatial\_field\_1d\_sim.py:70--75} and~\code{:98--107}). The spatial
pipeline chapters describe the matching theory side; here it is enough to know the
simulator realises the \emph{lattice} dispersion, a point we return to in the
caveats.

\subsection{Semi-implicit Crank--Nicolson (coupled reaction--diffusion)}

The coupled-RD simulator (\file{coupled\_rd\_1d\_sim.py}) integrates an
$N$-species reaction--diffusion system with cross-noise and unequal diffusion
constants. Its default scheme is \term{semi-implicit Crank--Nicolson}: a
trapezoidal-rule implicit step that, in this linear-stochastic setting, gives the
stationary covariance \emph{exactly to} $O(\dd t^{2})$ via a Pad\'e(1,1)
propagator (\file{coupled\_rd\_1d\_sim.py:30--44}). This is the most accurate of
the three improvements for the linear part. We return to this simulator, and its
\emph{analytic} oracle, in \S\ref{sec:coupled}.

\section{Correlated and multi-field noise (Cholesky)}

The two-field correlated simulators
(\code{sim\_ou\_quartic\_two\_dim\_corr\_numba} and
\code{sim\_ou\_quartic\_two\_dim\_color\_corr\_numba}) drive two fields with
jointly Gaussian white noise of correlation $\rho$:
\[
  \avg{\eta_\alpha(t)\,\eta_\beta(t')} \;=\; C_{\alpha\beta}\,\delta(t-t'),
  \qquad
  C \;=\; \begin{pmatrix} 1 & \rho \\ \rho & 1 \end{pmatrix}.
\]
Correlated Gaussians are produced by \term{Cholesky factorization} of the
covariance matrix $C$: writing $C = L L^{\mathsf T}$ with $L$ lower-triangular,
multiplying a vector of i.i.d.\ unit Gaussians by $L$ yields a Gaussian vector
with covariance exactly $C$. For the $2\times 2$ case the Cholesky factor is
\[
  L \;=\; \begin{pmatrix} 1 & 0 \\ \rho & \sqrt{1-\rho^{2}} \end{pmatrix},
\]
so the correlated drivers are (\file{ou\_langevin\_sim\_numba.py:351--352}):

\begin{lstlisting}
u = np.random.randn()
v = np.random.randn()
eta_x = u
eta_y = rho * u + rho_perp * v   # rho_perp = sqrt(1 - rho^2)
\end{lstlisting}

The correlation must satisfy $\rho \in [-1, 1]$ strictly; outside that range $C$
loses positive-semidefiniteness and no real Cholesky factor exists. The code
clamps $1-\rho^{2}$ to $\ge 0$ defensively
(\file{ou\_langevin\_sim\_numba.py:332--335}):

\begin{lstlisting}
one_minus_rho_sq = 1.0 - rho * rho
if one_minus_rho_sq < 0.0:
    one_minus_rho_sq = 0.0
rho_perp = np.sqrt(one_minus_rho_sq)
\end{lstlisting}

\begin{gotcha}
\textbf{An out-of-range $\rho$ yields NaN, not an error.} The clamp keeps the
square root from raising, but an out-of-range correlation then silently propagates
\code{NaN} through the trajectory rather than crashing
(\file{ou\_langevin\_sim\_numba.py:332--335}). If a two-field comparison plot comes
back empty, an out-of-bounds $\rho$ is the first thing to check.
\end{gotcha}

\section{The numba/Sage boundary}
\label{sec:numbasage}

This section is the single most important tooling fact in the subsystem. The
simulators are decorated with \code{numba}, and they live in plain \code{.py}
files for one specific reason that interacts with the SageMath kernel. To explain
it we first explain numba.

\subsection{What numba is}

\term{Numba} is a just-in-time (JIT) compiler for a subset of Python and NumPy. The
decorator \code{@numba.njit} (``no-python JIT'') compiles the decorated function to
native machine code, via the LLVM compiler infrastructure, the \emph{first} time
the function is called, then caches that compiled version for all subsequent
calls. The \code{njit} mode (also written \code{nopython=True}) forbids any
fallback to the slow Python interpreter: if numba cannot infer a concrete machine
type for every variable, it raises a compilation error rather than silently
running slowly. For the tight scalar arithmetic in an Euler step --- a loop of
millions of iterations of $x = x + \dd t\cdot\text{drift} + \sqrt{\dd t}\cdot
\text{noise}$ --- this buys roughly a $100\times$ speedup over pure Python (the
Hawkes header quotes ``\textasciitilde 15M steps/sec on Apple M1'',
\file{hawkes\_sim\_numba.py:37}).

Every Euler-step simulator is one \code{@numba.njit} function with a plain
scalar/array signature (\file{ou\_langevin\_sim\_numba.py:17--21}):

\begin{lstlisting}
@numba.njit
def sim_ou_quartic_numba(n_steps, dt_sim, mu, eps, D, x_init,
                         bin_size_steps, n_bins, seed):
\end{lstlisting}

\subsection{Why these must be \texttt{.py} files: the Sage-Integer trap}

When a notebook runs under the \term{SageMath kernel}, Sage's \term{preparser}
rewrites the source of every notebook cell before executing it. An integer literal
\code{0} becomes \code{Integer(0)} (a SageMath ring element), and \code{1.5}
becomes \code{RealNumber('1.5')}. These are \emph{not} Python \code{int}/\code{float};
they are objects from Sage's algebra system, carrying exact-arithmetic semantics.
Numba's type inference cannot handle them --- it only knows native machine types.
The docstring of the Hawkes simulator states the rule directly
(\file{hawkes\_sim\_numba.py:5--11}):

\begin{quote}\itshape
This file MUST be a plain \code{.py} file (not executed through SageMath's
preparser) because SageMath's preparser converts integer literals like \code{0}
to \code{Integer(0)}\,\ldots which Numba's type inference cannot handle.
\end{quote}

Two disciplines follow, and they are enforced throughout the subsystem:

\begin{enumerate}
  \item \textbf{The simulators live in \code{.py} modules, never in notebook
        cells.} Only \emph{cell source} is preparsed; imported \code{.py} files
        are read as ordinary Python. So \code{from models.ou\_langevin\_sim\_numba
        import sim\_ou\_quartic\_numba} is safe --- the function body never passes
        through the preparser.
  \item \textbf{Every argument at the call site must be cast to a plain Python
        \code{int}/\code{float}} before the \code{njit} function sees it. This is
        why the notebooks read \code{mu = float(fp['mu'])} rather than passing the
        resolved parameter straight through. The non-numba spatial simulators
        defend the same way \emph{inside} the function body --- \code{N = int(N);
        n\_steps = int(n\_steps); ...} (\file{spatial\_field\_1d\_sim.py:141--145})
        --- because \code{np.random.default\_rng} also rejects a Sage
        \code{Integer}.
\end{enumerate}

\begin{gotcha}
\textbf{An uncast Sage scalar crashes numba.} This is not stylistic. If you pass a
Sage \code{Integer} or \code{RealNumber} into an \code{@numba.njit} function you
get a numba typing error at \emph{compile} time; if you pass one into
\code{np.random.default\_rng} (the spatial path) you get a runtime rejection.
Reading parameters off \code{res['\_resolved']['parameters']} and casting with
\code{float(...)} / \code{int(...)} also guarantees the simulation uses
\emph{exactly} the values the theory used.
\end{gotcha}

\subsection{The two random-number APIs}

A small but real difference between the numba and non-numba simulators is which
NumPy random API they use. The numba-jitted Langevin and Hawkes simulators use the
\emph{legacy global} generator: \code{np.random.seed(seed)}
(\file{ou\_langevin\_sim\_numba.py:68}) seeds it, then \code{np.random.randn()}
(\file{ou\_langevin\_sim\_numba.py:82}) draws a single $N(0,1)$ and
\code{np.random.poisson(lam)} draws spike counts. The numba runtime supports this
legacy interface but not the modern generator object. The spatial simulators, which
are \emph{not} jitted, use the modern \code{np.random.default\_rng(seed)} generator
(\file{spatial\_field\_1d\_sim.py:154}) with \code{rng.standard\_normal(M)}, because
they are free to use the better API --- and, as noted above, that very API is what
forces the explicit \code{int(...)} casts inside their function bodies.

\subsection{JIT warmup}

Because the first call to an \code{njit} function pays the one-time compilation
cost, the notebooks make a throwaway warmup call before the timed measurement
loop, so the compilation is not counted in the timing:

\begin{lstlisting}
_ = sim_ou_quartic_numba(1000, dt_sim, mu, eps, D, 0.0,
                         bin_size_steps, 100, 0)   # JIT warmup
\end{lstlisting}

The arguments are throwaway (1000 steps, 100 bins, seed 0); the only purpose is to
trigger compilation. Every subsequent call in the run loop uses the cached native
code.

\subsection{What this subsystem does \emph{not} use}

For calibration against the rest of the manual: this subsystem uses \emph{no}
SageMath (it only avoids it), \emph{no} \code{nauty}/\code{networkx} (there is no
graph enumeration --- there are no diagrams here), and \emph{no} \code{sympy}
(there is no symbolic algebra --- the equations of motion are hand-coded as numeric
arithmetic). The two primary files touch only \code{numpy}, \code{numba}, and the
standard-library \code{collections.Counter}; the wider \file{models/} directory
additionally uses \code{scipy.linalg} in exactly one place
(\S\ref{sec:coupled}). That cleanness is the point.

\section{The cumulant estimator: factorial moments}
\label{sec:factorial}

We now turn to \file{models/cumulant\_estimator.py}, which converts a binned
trajectory into an empirical cumulant slice. The mathematical heart of the file is
the \emph{factorial moment}, and to motivate it we must understand a problem that
arises only for discrete point-process (spike-train) data.

\subsection{The shot-noise diagonal}

For point-process data, the raw $k$-point correlation contains delta-function
\term{contact terms}: a spike is perfectly correlated with itself. Concretely, the
same-bin product of counts $n^{2}$ overcounts --- it includes ``a spike paired with
itself,'' a shot-noise term scaling like $1/\dd t_{\mathrm{bin}}$ that is \emph{not}
part of the smooth connected correlator the diagrams compute. If you naively
estimated $\avg{n^{2}}$ you would measure the true covariance \emph{plus} this
spurious diagonal spike.

The fix is the \term{falling factorial}. Replace the ordinary power $n^{m}$ (at $m$
coincident same-population fields in the same bin) with
\begin{equation}
  n^{(m)} \;=\; n\,(n-1)\,(n-2)\cdots(n-m+1).
  \label{eq:falling}
\end{equation}
For a Poisson variable the identity $\avg{n^{(m)}} = \lambda^{m}$ holds
\emph{exactly} --- the falling factorial strips off precisely the self-coincidence
terms. The resulting estimator targets the \term{factorial cumulant density}: the
smooth, off-diagonal part of the connected $k$-point cumulant, with the
$\delta$-function diagonals removed (\file{cumulant\_estimator.py:8--31}). For $k=2$
this reduces to the ordinary rate covariance, since the only diagonal is the shot
noise $\delta(\tau)$ (\file{cumulant\_estimator.py:13--16}). The reference is Daley
\& Vere-Jones, \emph{An Introduction to the Theory of Point Processes}, Ch.~5 on
factorial moment measures (\file{cumulant\_estimator.py:55--57}).

The falling factorial is built elementwise over a whole bin array
(\file{cumulant\_estimator.py:63--81}):

\begin{lstlisting}
def falling_factorial_array(n_arr, m):
    result = np.ones_like(n_arr, dtype=float)
    for j in range(m):
        result *= (n_arr - j)
    return result
\end{lstlisting}

For $m=1$ this is just $n$ itself (linear centering, no correction); for $m \ge 2$
it removes the self-coincidence shot noise.

\subsection{Rate conversion}

A spike-counting field's value $n$ is a \emph{count} per bin. To convert it to a
rate (units of Hz) it is divided by the bin width, $n/\dd t_{\mathrm{bin}}$. This is
the \code{denom = dt\_bin} that appears throughout the estimator. Smooth voltage
fields (field type \code{'dv'}) are \emph{not} counts and use \code{denom = 1.0}
(no rate conversion; units of volts). The choice is made by a small closure
(\file{cumulant\_estimator.py:218--221}):

\begin{lstlisting}
def _centering_denom(leg_idx):
    return dt_bin if _is_spike_ft(field_types[leg_idx]) else 1.0
\end{lstlisting}

\section{The estimator algorithm, step by step}

The core routine is \code{compute\_kpoint\_slice}
(\file{cumulant\_estimator.py:84}). Its signature exposes the full machinery:

\begin{lstlisting}
def compute_kpoint_slice(binned_counts, dt_bin, pop_indices, lag_bins,
                         max_lag_bins, n_fft=None,
                         field_types=None, voltage_bins=None,
                         ext_binned_counts=None):
\end{lstlisting}

The idea is: hold all but one of the $k$ external legs at fixed time lags, sweep
the remaining leg's lag across a window, and report the connected cumulant as a
function of that one sweep lag. The leg that sweeps is the \term{sweep axis}; the
others are pinned. We walk the ten steps in order.

\paragraph{1. Validate the sweep axis} (\file{cumulant\_estimator.py:139--145}).
Exactly one entry of \code{lag\_bins} must be \code{None}; that index is recorded
as \code{sweep\_idx}. An assertion enforces this.

\paragraph{2. Normalize field types} (\file{cumulant\_estimator.py:148--202}).
Each leg has a field type; the default is \code{'dn'}. The routine raises if a
\code{'dv'} leg is requested without \code{voltage\_bins}, or a \code{'dm'} leg
without \code{ext\_binned\_counts}, and maps natural names through the nested
\code{\_normalize\_ft} helper. We discuss the field-type system and the
bug \code{\_normalize\_ft} fixed in \S\ref{sec:fieldtypes}.

\paragraph{3. Define helpers} (\file{cumulant\_estimator.py:204--221}).
\code{\_data\_for(leg)} picks the counts/voltage/external array by field type;
\code{\_is\_spike\_ft} is \code{True} for \code{'dn'}/\code{'dm'}; and
\code{\_centering\_denom} returns \code{dt\_bin} for spike fields, \code{1.0} for
voltage (shown above).

\paragraph{4. Determine the valid window} (\file{cumulant\_estimator.py:225--240}).
From the fixed lags (plus an implicit reference at lag 0) the routine computes
\code{valid\_start = max(0, -min\_lag)}, \code{valid\_end = n\_bins - max(0,
max\_lag)}, and \code{n\_valid = valid\_end - valid\_start}. This is the set of
reference bins for which \emph{every} fixed leg still lands inside the array --- no
wraparound. It raises if \code{n\_valid <= 2*max\_lag\_bins}, i.e.\ if the
recording is too short for the requested lag window.

\paragraph{5. Means} (\file{cumulant\_estimator.py:242--251}). The mean of each
\code{(pop, field\_type)} combination is computed \emph{on the valid window}, for
consistency with the product that follows. These means are what we subtract to
center each field.

\paragraph{6. Build the fixed-leg product} (\file{cumulant\_estimator.py:261--300}).
The fixed legs are grouped by lag value; within each lag, legs are grouped by
\code{(pop, ft)} via a \code{Counter} (\S\ref{sec:counter}). For multiplicity
\code{m == 1} (or any voltage leg) the field is centered linearly,
$(n - \text{mean})/\text{denom}$. For spike fields with \code{m >= 2} the falling
factorial is applied and then centered:

\begin{lstlisting}
if m == 1 or ft == 'dv':
    factor = (n_arr - mean_n) / denom
    if m > 1:
        factor = factor ** m          # voltage power, no factorial
else:
    raw = falling_factorial_array(n_arr, m) / denom**m
    factor = raw - raw.mean()         # spike: factorial, then centered
product *= factor
\end{lstlisting}

Note the asymmetry: a \emph{voltage} field at multiplicity $m>1$ is raised to the
power $m$ (it is smooth, with no shot noise to remove), whereas a \emph{spike}
field at $m \ge 2$ gets the factorial correction.

\paragraph{7. Build the sweep series} (\file{cumulant\_estimator.py:302--307}).
The swept field is centered the same way: $(\text{sweep} - \text{mean})/
\text{denom}$.

\paragraph{8. Cross-correlate by zero-padded FFT}
(\file{cumulant\_estimator.py:309--320}). The sweep-axis cross-correlation
$C(L) = \sum_t \text{product}(t)\cdot\text{sweep}(t+L)$ is computed with a
zero-padded FFT, using the convolution theorem (a cross-correlation is an inverse
FFT of one transform times the conjugate of the other):

\begin{lstlisting}
F_product = np.fft.fft(product, n=n_fft)
F_sweep   = np.fft.fft(sweep_rate, n=n_fft)
raw_xcorr = np.fft.ifft(F_product * np.conj(F_sweep)).real
\end{lstlisting}

The FFT must be zero-padded to \code{n\_fft >= 2*n\_valid} (the default is the next
power of two $\ge 2\,n_{\mathrm{valid}}$, \file{cumulant\_estimator.py:310--311})
so that the \emph{circular} convolution the FFT computes coincides with the desired
\emph{linear} correlation over the valid window, with no periodic-wraparound
contamination.

\begin{gotcha}
\textbf{The FFT lag indexing is fragile.} The cross-correlation lag is read as
\code{raw\_xcorr[(-lag) \% n\_fft]}. The comment in the code is load-bearing
(\file{cumulant\_estimator.py:316--319}): \code{ifft(F\_product * conj(F\_sweep))[L]}
equals $\sum_t \text{product}(t)\cdot\text{sweep}(t-L)$, so the value we want at
$+\text{lag}$ lives at $L = -\text{lag}$. The comment warns explicitly:
\textbf{``Do NOT use \texttt{[::-1]} --- that shifts the index by 1 bin.''} A
reversed-array shortcut introduces a one-bin offset that is easy to miss and
poisons the whole slice.
\end{gotcha}

\paragraph{9. Extract and normalize} (\file{cumulant\_estimator.py:322--332}). For
each lag in $[-\code{max\_lag\_bins}, +\code{max\_lag\_bins}]$ the raw correlation
is divided by the number of \emph{overlapping} valid bin pairs,
\code{n\_overlap = n\_valid - abs(lag)}, not by a constant. This is the
\term{lag-dependent normalization} that makes the finite-window estimator
unbiased: fewer pairs contribute at larger lags, and dividing by the actual count
corrects for it.

\paragraph{10. The $k=4$ cumulant subtraction} (\file{cumulant\_estimator.py:346--433}).
Discussed next.

\subsection{$k=2$, $k=3$, $k=4$, and beyond}

For $k \le 3$ the connected cumulant \emph{equals} the centered moment, so the
centered, factorial-corrected product computed above \emph{is} the cumulant ---
there is nothing further to subtract. This is why the comment at
\file{cumulant\_estimator.py:343--345} notes that for $k \le 3$ the cumulant equals
the central moment.

For $k=4$ the connected cumulant requires subtracting the three disconnected pair
products (\file{cumulant\_estimator.py:334--359}):
\begin{equation}
  \kappa_4(X_1,X_2,X_3,X_4)
  \;=\; m_4
       \;-\; m_2(1,2)\,m_2(3,4)
       \;-\; m_2(1,3)\,m_2(2,4)
       \;-\; m_2(1,4)\,m_2(2,3),
  \label{eq:kappa4}
\end{equation}
with $m_n$ the centered moments. The cross-correlation above gave $m_4$ at each
sweep lag; the three pair products are computed by a nested helper
\code{\_two\_point(leg\_a, leg\_b, lag\_value)}
(\file{cumulant\_estimator.py:361--420}). That helper is itself
factorial-corrected when the pair is same-population, same spike field type, and
same effective lag (a genuine shot-noise coincidence to remove), and uses ordinary
linear covariance in every other case (\file{cumulant\_estimator.py:381--420}).
The three partitions of four legs into two pairs are
$\{(0,1)(2,3),\,(0,2)(1,3),\,(0,3)(1,2)\}$
(\file{cumulant\_estimator.py:357--359}), and the disconnected piece is subtracted
lag by lag (\file{cumulant\_estimator.py:422--433}).

\begin{gotcha}
\textbf{$k>4$ is not a cumulant.} For $k>4$ the partition-based subtraction is
\emph{not} implemented. The estimator emits a \code{RuntimeWarning} and returns the
centered $k$-point \emph{moment} instead of the connected cumulant
(\file{cumulant\_estimator.py:435--440}). Separately, the full diagonal/contact
$\delta$-terms are \emph{never} estimated by this code --- they are distributions,
not functions on a $\tau$ grid (\file{cumulant\_estimator.py:33--40}). The
estimator targets the smooth off-diagonal part only.
\end{gotcha}

\subsection{\texttt{collections.Counter}: counting coincident legs}
\label{sec:counter}

\code{collections.Counter} is a standard-library \code{dict} subclass that counts
hashable items: \code{Counter(['a','a','b'])} returns \code{\{'a': 2, 'b': 1\}}.
The estimator uses it to find the \emph{multiplicity} $m$ --- how many legs of the
same \code{(population, field\_type)} land in the same bin, which is exactly the
order of the falling factorial to apply (\file{cumulant\_estimator.py:279--282}):

\begin{lstlisting}
pop_ft_multiplicities = Counter(
    (pop_indices[i], field_types[i]) for i in leg_indices)
for (pop, ft), m in pop_ft_multiplicities.items():
    ...
\end{lstlisting}

The same idiom appears in the direct estimator (\file{cumulant\_estimator.py:547}).
The value $m$ then selects between linear centering ($m=1$) and factorial
correction ($m \ge 2$).

\section{Field types: \texttt{dn}, \texttt{dv}, \texttt{dm}}
\label{sec:fieldtypes}

The estimator's \code{field\_types} argument labels each leg as one of three kinds
(\file{cumulant\_estimator.py:109--121}):

\begin{description}
  \item[\code{'dn'}] a cortical spike-train fluctuation: discrete counts in
        \code{binned\_counts}, shot-noise corrected with the falling factorial,
        rate-converted by $\dd t_{\mathrm{bin}}$.
  \item[\code{'dv'}] a voltage fluctuation: a smooth continuous field in
        \code{voltage\_bins}, linearly centered, no factorial, no rate conversion.
        \emph{This is the label the OU simulators use} --- the OU variable $x$ is a
        smooth field, so it is passed as \code{voltage\_bins} with
        \code{field\_types=['dv', ...]}.
  \item[\code{'dm'}] an external GTaS rate fluctuation: spike-like, in
        \code{ext\_binned\_counts}, treated exactly like \code{'dn'} (discrete
        counts, $\dd t_{\mathrm{bin}}$ denominator, factorial correction).
\end{description}

\begin{gotcha}
\textbf{The field-type normalization fixed a $1/\dd t_{\mathrm{bin}}^{2}$ bug.}
Before the \code{\_normalize\_ft} helper (\file{cumulant\_estimator.py:166--200}),
natural-name legs like \code{'n'} or \code{'nE'} were treated as non-spike, so they
used \code{denom = 1.0} instead of \code{dt\_bin} and the estimator returned
$\text{count}^{2}$ instead of $\text{rate}^{2}$ --- wrong by a factor
$1/\dd t_{\mathrm{bin}}^{2}$ per leg (a factor of 16 at $\dd t_{\mathrm{bin}} =
0.25$). Every notebook that built \code{field\_types = [ef[0] for ef in
external\_fields]} hit it. New code should still prefer the canonical
\code{'dn'}/\code{'dv'}/\code{'dm'} labels; the normalizer is a safety net, not a
license to be sloppy.
\end{gotcha}

\section{The convenience wrapper: \texttt{estimate\_kpoint\_slices}}

Notebooks rarely call \code{compute\_kpoint\_slice} directly. They call
\code{estimate\_kpoint\_slices} (\file{cumulant\_estimator.py:445}), which produces
all $k-1$ axis-parallel slices of the connected $k$-point cumulant in one go,
laid out to line up panel-for-panel with the theory side's
\code{res['C\_tau\_slices']}:

\begin{lstlisting}
def estimate_kpoint_slices(dt_bin, pop_indices, field_types, base_bins,
                           max_lag_bins, binned_counts=None, voltage_bins=None,
                           ext_binned_counts=None):
\end{lstlisting}

Leg 0 is the \emph{anchor}, pinned at lag 0. \code{base\_bins} is a length-$(k-1)$
list of integer bin offsets --- the ``base point,'' one entry per non-anchor leg.
For each $j = 1\ldots k-1$ the wrapper builds a lag list with leg 0 at 0, every
other non-anchor leg at its \code{base\_bins} offset, and leg $j$ set to
\code{None} (the sweep), then calls \code{compute\_kpoint\_slice} and stacks the
result (\file{cumulant\_estimator.py:472--483}):

\begin{lstlisting}
for j in range(1, k):
    lag = [0] * k                       # leg 0 anchored at 0
    for i in range(1, k):
        lag[i] = int(base_bins[i - 1])  # non-anchor legs at the base...
    lag[j] = None                       # ...except leg j, which sweeps
    tau, C = compute_kpoint_slice(...)
    rows.append(np.asarray(C))
\end{lstlisting}

It returns \code{(tau\_grid, C)} with \code{tau\_grid} of shape
\code{(2*max\_lag\_bins+1,)} and \code{C} of shape
\code{(k-1, 2*max\_lag\_bins+1)} --- row $j-1$ is slice $j$. As a convenience, if
\code{binned\_counts} is \code{None} but \code{voltage\_bins} is given, it
allocates a zero \code{binned\_counts} (the pure-voltage case used by the OU
simulators, \file{cumulant\_estimator.py:469--470}). For $k=2$ this returns a
single autocorrelation slice; for $k\ge 3$ it returns one row per swept leg, which
is what makes the simulation curves line up with \code{dd.plot\_cumulant}'s $k\ge3$
slice panels.

\section{The non-FFT cross-check: \texttt{compute\_kpoint\_slice\_direct}}

\code{compute\_kpoint\_slice\_direct} (\file{cumulant\_estimator.py:486}) has the
same interface and output as \code{compute\_kpoint\_slice} but uses an explicit
double loop over \code{(sweep\_lag, t)} instead of FFTs --- $O(n_{\mathrm{bins}}
\times n_{\mathrm{lags}})$ versus $O(n_{\mathrm{bins}}\log n_{\mathrm{bins}})$. It is
slower but free of any FFT artifact (circular wrap, zero-pad, spectral leakage),
which makes it the natural way to cross-check the FFT estimator.

It is, deliberately, a \emph{reduced} re-implementation. It handles spike counts
only --- there are no \code{field\_types}, \code{voltage\_bins}, or
\code{ext\_binned\_counts} parameters --- and its $m \ge 2$ factorial mean is
approximated by the \emph{full-array} mean rather than the valid-window mean used
by the FFT path (\file{cumulant\_estimator.py:559--563}):

\begin{lstlisting}
ff_mean = falling_factorial_array(
    binned_counts[pop], m
).mean() / dt_bin**m
val *= (ff - ff_mean)
\end{lstlisting}

\begin{gotcha}
\textbf{Three estimators, one logic, no shared kernel.}
\code{compute\_kpoint\_slice} (FFT), \code{compute\_kpoint\_slice\_direct} (non-FFT),
and \code{cumulant\_direct\_numba.\_kpoint\_slice\_direct\_k3} (a numba-compiled
$k=3$ variant) implement the \emph{same} factorial-cumulant logic three times for
three different reasons --- FFT speed, FFT-free validation, and numba-compiled $k=3$
throughput. They must be kept in sync by discipline; there is no shared kernel.
Small discrepancies between the FFT path and the direct path are \emph{expected} at
the window edges, because the direct path uses the full-array factorial mean rather
than the valid-window mean.
\end{gotcha}

\section{Spatial estimators}

The spatial simulators (\file{spatial\_field\_1d\_sim.py} and its 2D/3D siblings)
carry their \emph{own} estimators, because a spatial field is averaged over the
periodic translation group rather than over a temporal lag sweep. The 1D file
exposes (\file{spatial\_field\_1d\_sim.py}):

\begin{itemize}
  \item \code{structure\_factor(snaps, meta)} --- the momentum-space two-point
        function $S(q) = \avg{|\hat\phi_q|^{2}}$ (\code{:175}), which tends to
        $T/(\mu + Dq^{2})$ at tree level.
  \item \code{equal\_time\_correlator(snaps)} --- the real-space equal-time
        correlator $C(x)$ (\code{:191}).
  \item \code{lattice\_sum\_variance(L, N, mu, D, T)} --- the exact $\lambda = 0$
        variance reference (\code{:205}), summed over the \emph{lattice} modes.
  \item \code{third\_cumulant\_x(snaps)} --- $\kappa_3(m) =
        \avg{\delta\phi(x_0)\,\delta\phi(x_0+m\,\dd x)^{2}}$ with a standard error
        (\code{:215}).
\end{itemize}

These are the spatial counterparts of \code{compute\_kpoint\_slice}: instead of a
temporal lag sweep they use the periodic translational average --- one circular
cross-correlation per snapshot via FFT.

\begin{gotcha}
\textbf{Matched-cutoff principle.} The spatial simulator realises the
\emph{lattice} dispersion $\omega = \mu + (2D/\dd x^{2})(1-\cos)$, not the continuum
$\mu + Dq^{2}$. Quantitative validation must therefore compare against a theory
computed at the \emph{same} UV cutoff $k_{\max} = \pi/\dd x$, or against the lattice
oracle (\code{lattice\_sum\_variance}, or
\code{coupled\_box\_correlator(dispersion='lattice')}). The two dispersions differ
at $O(\dd x)$ --- about 4\% on $C(0,0)$ at $L=20$, $n_x=128$.
\end{gotcha}

\section{Coupled multi-field validation}
\label{sec:coupled}

The coupled reaction--diffusion simulator
(\code{simulate\_coupled\_rd\_1d}, \file{coupled\_rd\_1d\_sim.py}) is the ground
truth for the coupled Dyson series. Uniquely in this subsystem, it ships with an
\emph{analytic} oracle as well as a simulator, and the oracle is the one place
\code{scipy} appears.

\term{SciPy} is the scientific-computing layer built on NumPy;
\code{scipy.linalg} is its dense linear-algebra module (a set of LAPACK wrappers).
The oracle \code{coupled\_box\_correlator} uses two of its routines, per Fourier
mode (\file{coupled\_rd\_1d\_sim.py:73}):

\begin{itemize}
  \item \code{scipy.linalg.expm(A)} --- the \term{matrix exponential} $\ee^{A}$,
        used to build the exact lag propagator $G(\tau) = \ee^{-A(q)\,|\tau|}$ at
        each mode (\file{coupled\_rd\_1d\_sim.py:353}).
  \item \code{scipy.linalg.solve\_continuous\_lyapunov(A, Q)} --- solves the
        \term{continuous Lyapunov equation} $A\Sigma + \Sigma A^{\mathsf T} = Q$ for
        $\Sigma$, the exact stationary covariance of the linear coupled system at
        each mode (\file{coupled\_rd\_1d\_sim.py:351}).
\end{itemize}

Together these give an \emph{analytic} (no-simulation) reference against which the
coupled Dyson series is validated. The simulator side
(\code{simulate\_coupled\_rd\_1d}) returns a dict
\code{\{'C', 'C\_err', 'C\_rep', 'x\_grid', 'taus', 'meta'\}}.

\begin{gotcha}
\textbf{Coupled-RD Courant guard.} \code{simulate\_coupled\_rd\_1d} \emph{asserts}
$\dd t\cdot\max(D)/\dd x^{2} < 0.4$ (\file{coupled\_rd\_1d\_sim.py:212--214}) --- an
\code{AssertionError}, not a graceful message --- to keep both schemes stable and
the Pad\'e lag bias negligible. If you crank the step or shrink the grid spacing
too far, this is what fires.
\end{gotcha}

\section{The rest of \texttt{models/}}

For completeness, the wider supporting cast. All are independent brute-force
validators paired with the estimator described above.

\begin{center}
\small
\begin{longtable}{@{}p{0.30\textwidth} p{0.15\textwidth} p{0.47\textwidth}@{}}
\toprule
\textbf{File} & \textbf{Family} & \textbf{What it simulates} \\
\midrule
\endhead
\file{cumulant\_direct\_numba.py} & estimator &
  \code{@numba.njit} direct (non-FFT) $k=3$ factorial-cumulant slice
  (\code{\_kpoint\_slice\_direct\_k3}); a plain \code{.py} file so numba can
  compile it. \\
\file{hawkes\_sim\_numba.py} & Hawkes &
  Single-population linear Hawkes; Poisson spikes feed the voltage; returns the
  tuple \code{(binned\_counts, voltage\_bins, total\_spikes)}. \\
\file{hawkes\_sim\_multipop\_numba.py} & Hawkes &
  Multipopulation linear Hawkes with per-pair exponential synaptic filters
  $F_{ij}$; flat E/I neuron index; matches
  \file{multipopulation\_test.theory.py}. \\
\file{hawkes\_sim\_expg\_numba.py}, \file{hawkes\_sim\_quad\_expg\_numba.py},
  \file{hawkes\_sim\_*\_gtas\_numba.py} & Hawkes &
  Exponential-$g$ synaptic-filter and quadratic-rate variants; \code{gtas} =
  ``GTaS'' external spike drive (the \code{'dm'} field, fourth return value
  \code{ext\_binned\_counts}). \\
\file{hawkes\_sage.py} and the \file{hawkes\_*\_sage} / \file{hawkes\_*\_expg}
  family & Hawkes (Sage) &
  Sage-side theory/orchestration helpers for the Hawkes models (these \emph{are}
  preparsed; not \code{njit}). \\
\file{dendritic\_linear\_sim\_numba.py} & Hawkes-like &
  Dendritic linear model simulator; imports \code{compute\_kpoint\_slice}. \\
\file{spatial\_field\_1d\_sim.py} & spatial Langevin &
  1D scalar $\phi^{4}$ on a ring, spectral ETD1; \code{simulate},
  \code{structure\_factor}, \code{equal\_time\_correlator},
  \code{lattice\_sum\_variance}, \code{third\_cumulant\_x}. \\
\file{spatial\_field\_phi6\_1d\_sim.py} & spatial Langevin &
  1D $\phi^{6}$ (cubic+quintic force); near-copy of the $\phi^{4}$ simulator. \\
\file{spatial\_field\_2d\_sim.py}, \file{spatial\_field\_3d\_sim.py} &
  spatial Langevin &
  2D scalar on a torus / 3D scalar, ETD1; radial structure-factor and correlator
  estimators. \\
\file{coupled\_rd\_1d\_sim.py} & coupled RD &
  $N$-species reaction--diffusion with cross-noise and unequal $D$; semi-implicit
  Crank--Nicolson or explicit scheme; \emph{plus} the analytic oracle
  \code{coupled\_box\_correlator} (\code{scipy.linalg.expm} /
  \code{solve\_continuous\_lyapunov}). \\
\bottomrule
\end{longtable}
\end{center}

\section{Putting it together: a sim-compare notebook}

We close with the canonical wiring, taken from cell~9 of
\file{notebooks/examples/temporal\_ou\_quartic\_white.ipynb} --- the notebook
behind the Chapter~\ref{ch:quickstart} payoff. Read it against the seven salient
points that follow:

\begin{lstlisting}
from models.ou_langevin_sim_numba import sim_ou_quartic_numba
from models.cumulant_estimator import estimate_kpoint_slices

fp = res['_resolved']['parameters']          # same physics as the theory run
mu = float(fp['mu']); eps = float(fp['eps']); D = float(fp['D'])   # CAST to float!

dt_sim, dt_bin = 0.01, 0.02
n_steps        = int(T_sim / dt_sim)
bin_size_steps = int(max(round(dt_bin / dt_sim), 1))
dt_bin_eff     = bin_size_steps * dt_sim     # the EFFECTIVE bin width (use this)
n_bins         = int(n_steps // bin_size_steps)
max_lag_bins   = int(tau_max / dt_bin_eff)

k           = int(res['_resolved']['k'])
base_bins   = [int(round(b / dt_bin_eff)) for b in base]   # length k-1
pop_indices = [0]*k
field_types = ['dv']*k                       # OU field x is a 'dv' (smooth)

_ = sim_ou_quartic_numba(1000, dt_sim, mu, eps, D, 0.0,
                         bin_size_steps, 100, 0)            # JIT warmup
for r in range(N_RUNS):
    x_bins = sim_ou_quartic_numba(n_steps, dt_sim, mu, eps, D, 0.0,
                                  bin_size_steps, n_bins, rng_base + r)
    tau_sim, Cj = estimate_kpoint_slices(
        dt_bin_eff, pop_indices, field_types, base_bins,
        max_lag_bins, voltage_bins=x_bins)                 # Cj: (k-1, n_tau)
    C_runs.append(np.asarray(Cj).real)
C_sim = np.array(C_runs).mean(axis=0)
C_err = np.array(C_runs).std(axis=0, ddof=1) / np.sqrt(N_RUNS)
\end{lstlisting}

\begin{enumerate}
  \item \textbf{The simulator is handed \code{res['\_resolved']['parameters']}} ---
        the pipeline-resolved physics --- so theory and simulation share parameters
        but \emph{not} code.
  \item \textbf{Every scalar is cast} (\code{float(fp['mu'])}, \code{int(...)})
        before the \code{njit} call: the Sage-Integer guard of \S\ref{sec:numbasage}
        in action.
  \item \textbf{A throwaway warmup call} triggers JIT compilation once so it is not
        timed in the measurement loop.
  \item \textbf{The OU field is passed as \code{voltage\_bins} with
        \code{field\_types=['dv']}} --- it is a smooth continuous field, so it uses
        linear centering, not the factorial spike correction. (Hawkes spike legs
        would use \code{'dn'} and pass \code{binned\_counts} instead.)
  \item \textbf{\code{dt\_bin\_eff}} (= \code{bin\_size\_steps}\,$\times$\,\code{dt\_sim}),
        not the nominal \code{dt\_bin}, is what the estimator receives --- the
        binning rounds \code{dt\_bin/dt\_sim} to an integer number of steps, so the
        true bin width may differ slightly from the requested one.
  \item \textbf{Multiple independent runs} (\code{N\_RUNS}, distinct seeds) give the
        error bar \code{C\_err = std/sqrt(N\_RUNS)}.
  \item \textbf{The comparison} is \code{C\_sim} versus \code{res['C\_tau\_by\_ell'][0]}
        (tree) and \code{res['C\_tau']} (loop-corrected) at the $\tau=0$ index ---
        \emph{that overlap is the trust-chain verdict.}
\end{enumerate}

For $k\ge 3$, \code{estimate\_kpoint\_slices} returns one row per swept leg, lining
up panel-for-panel with \code{dd.plot\_cumulant}'s theory slices, so the same loop
validates a third or fourth cumulant with no change beyond \code{k} and
\code{base\_bins}.

\section{Caveats checklist}

A consolidated list of the traps documented above, for quick reference when a
comparison plot misbehaves:

\begin{enumerate}
  \item \textbf{Sage \code{Integer}/\code{RealNumber} crash numba.} Keep the
        simulators in \code{.py} files; cast every call-site scalar with
        \code{int()}/\code{float()} (\S\ref{sec:numbasage}).
  \item \textbf{Colored-noise factor of 2.} The white limit of
        \code{sim\_ou\_quartic\_colored\_numba} is $\avg{\xi\xi}\to 4D\,\delta$;
        use $D_{\mathrm{white}} = 2\,D_{\mathrm{colored}}$ to compare
        (\S\ref{sec:colored}).
  \item \textbf{\code{sim\_ou\_quartic\_two\_dim\_corr\_numba} vs its theory file
        disagree on the noise coefficient.} The simulator uses
        $\avg{\xi\xi} = 2D\,\delta$; the theory file as written declares
        \code{coefficient='D1'} ($= D\,\delta$, half strength). A quantitative
        compare needs one side fixed --- edit the theory to \code{'2*D1'} or halve
        the simulator's $D$ (\file{ou\_langevin\_sim\_numba.py:412--419}). Speed
        tests are unaffected.
  \item \textbf{$\rho$ out of $[-1,1]$ yields NaN, not an error}
        (\file{ou\_langevin\_sim\_numba.py:332--335}).
  \item \textbf{FFT lag indexing is fragile.} Read as
        \code{raw\_xcorr[(-lag) \% n\_fft]}; do not use \code{[::-1]}; zero-pad to
        \code{n\_fft >= 2*n\_valid} (\file{cumulant\_estimator.py:316--319}).
  \item \textbf{Field-type normalization fixed a $1/\dd t_{\mathrm{bin}}^{2}$ bug.}
        Prefer the canonical \code{'dn'}/\code{'dv'}/\code{'dm'} labels
        (\S\ref{sec:fieldtypes}).
  \item \textbf{$k>4$ is not a cumulant} --- the estimator returns the centered
        moment with a \code{RuntimeWarning}
        (\file{cumulant\_estimator.py:435--440}).
  \item \textbf{Euler--Maruyama bias.} All \code{*\_numba} Langevin simulators are
        weak order 1; stationary stats carry $O(\dd t)$ bias, so
        $\dd t_{\mathrm{sim}}\ll 1/\mu$ is required. The colored and spatial
        simulators remove this for the \emph{linear} part via the exact OU / ETD1
        update; the coupled-RD semi-implicit scheme is $O(\dd t^{2})$-exact in the
        stationary covariance.
  \item \textbf{Coupled-RD Courant guard.} \code{simulate\_coupled\_rd\_1d}
        asserts $\dd t\cdot\max(D)/\dd x^{2} < 0.4$
        (\file{coupled\_rd\_1d\_sim.py:212--214}).
  \item \textbf{Bin finely or bias $C_{xx}(0)$ low.} A coarse bin smooths the field
        and underestimates the equal-time variance; the example uses
        \code{dt\_bin}\,$=0.02$ for $\mu\sim1$.
  \item \textbf{Matched-cutoff principle (spatial).} Compare against a theory at the
        same $k_{\max}=\pi/\dd x$, or against the lattice oracle; the dispersions
        differ at $O(\dd x)$.
\end{enumerate}
